% Introduction: EMA, MACD/SIGNAL equations, trading algorithm.

\begin{frame}
  \frametitle{EMA wygładza ceny, nadając największą wagę ostatnim notowaniom}
  \begin{equation*}
    \alpha_N = \frac{2}{N+1},\qquad
    E_0^{(N)} = x_0,\qquad
    E_i^{(N)} = \alpha_N x_i + (1 - \alpha_N)\,E_{i-1}^{(N)}
    \quad\text{dla } i \ge 1
    \tag{1}
  \end{equation*}
  \vskip1mm
  \begin{itemize}
    \item $x_i$ --- cena zamknięcia w próbce $i$
    \item $N$ --- okres średniej
    \item $\alpha_N$ --- współczynnik wygładzania
    \item $E_i^{(N)}$ --- wartość EMA o okresie $N$
  \end{itemize}
  \vskip1mm
  Krótszy okres $N$ oznacza szybszą reakcję na zmiany ceny.
\end{frame}

\begin{frame}
  \frametitle{MACD mierzy rozbieżność dwóch średnich EMA, a SIGNAL ją wygładza}
  MACD (ang. moving average convergence/divergence) to wskaźnik analizy
  technicznej, który na podstawie cen zamknięcia pomaga ocenić kierunek
  i dynamikę zmian cen oraz wyznacza sygnały kupna i sprzedaży.
  \begin{equation*}
    M_i = E_i^{(N_1)} - E_i^{(N_2)}
    \tag{2}
  \end{equation*}
  \begin{equation*}
    S_0 = M_0,\qquad
    S_i = \alpha_{N_3} M_i + (1 - \alpha_{N_3})\,S_{i-1}
    \quad\text{dla } i \ge 1
    \tag{3}
  \end{equation*}
  \vskip2mm
  \begin{itemize}
    \item $M$ --- linia MACD
    \item $S$ --- linia SIGNAL
    \item $N_1$ i $N_2$ --- okresy krótszej i dłuższej EMA cen,
      $N_3$ --- okres wygładzania SIGNAL
    \item $\alpha_{N_3} = 2/(N_3+1)$ jak we wzorze (1);
      w projekcie $(N_1, N_2, N_3) = (12, 26, 9)$
  \end{itemize}
\end{frame}

\begin{frame}
  \frametitle{Zmiana znaku histogramu wyznacza momenty kupna i sprzedaży}
  \begin{equation*}
    H_i = M_i - S_i
    \tag{4}
  \end{equation*}
  \vskip2mm
  \begin{itemize}
    \item Kupno w próbce $i$, gdy $H_{i-1} \le 0$ i $H_i > 0$.
    \item Sprzedaż, gdy $H_{i-1} \ge 0$ i $H_i < 0$; bez interpolacji czasu
      przecięcia.
    \item Symulacja: sygnały od próbki 101, realizacja kolejnej sesji
      ($d = 1$), pełne zaangażowanie kapitału.
  \end{itemize}
\end{frame}
